\documentclass[10pt,a4paper]{amsart}
\usepackage[margin=19mm]{geometry}
\usepackage{amsmath,amssymb,mathtools}
\usepackage[scr=rsfs,cal=euler]{mathalfa}
\usepackage{xfrac}
\usepackage[unicode,hidelinks,bookmarks=false]{hyperref}
\newcommand{\ie}{i.e.\ }
\newcommand{\cf}{cf.\ }
\newcommand{\resp}{respectively\ }
\newcommand{\Subsection}{\subsection}
\providecommand{\nobreakdash}{\nobreak\hbox{-}}
\setlength{\emergencystretch}{2em}
\multlinegap0pt
\usepackage{amssymb}
\usepackage{mathtools}
\newtheorem{definition}{Definition}
\newtheorem{theorem}{Theorem}
\newcommand{\C}{\mathscr{C}}
\newcommand{\E}{\mathfrak{E}}
\newcommand{\Fix}{\operatorname{Fix}}
\def\R{\mathbb{R}}
\title{Perturbation-resilient inertial Krasnosel'skii-type hybrid retractions for generalized nonexpansive mappings}
\author{Markjoe O. Uba}
\date{Source: arXiv:2606.00528v3 (CC BY 4.0)}
\begin{document}
\maketitle

\noindent\emph{Source and licence.} Perturbation-resilient inertial Krasnosel'skii-type hybrid retractions for generalized nonexpansive mappings. Version arXiv:2606.00528v3 (CC BY 4.0), distributed by arXiv under the Creative Commons Attribution 4.0 International licence, http://creativecommons.org/licenses/by/4.0/, as recorded in the arXiv licence metadata for this exact version and shown on the abstract page https://arxiv.org/abs/2606.00528v3. Full licence notice: \texttt{licenses/arxiv-2606.00528-CC-BY-4.0.txt}.\par\smallskip

\noindent\textit{Editorial scope.} Counted unit: the article's theorem thm:bregman (source0.tex lines 515--537). The article's complete original proof of it is delivered with it (source0.tex lines 539--658, one \texttt{proof} environment of 120 lines). No mathematics is written, repaired or re-derived: the statement and the proof are the article's own lines, in the article's own notation, and no retained line is substituted or re-flowed.  Retained uncounted as the source-local context the proof needs: lines 432--444; lines 477--481; lines 494--496; lines 502--505.  Nothing was omitted from any retained span, so the package carries no omission record.  No retained line contains a citation, so the wrapper carries no bibliography and no bibliography entry.  No retained line contains a figure environment, an image-inclusion command or drawing code, so no image is delivered and the package declares no asset dependency.  Editorial formatting only, in addition: an AMS \texttt{amsart} preamble that restates the article's own declarations and macros used by the retained lines, namely the environments \texttt{definition}, \texttt{theorem} on separate counters, together with the standard packages those lines need.  The retained environments print in this document as Definition 1, Theorem 1 in order of appearance, whereas the article prints the same items with the numbering of its own full document.  The complete native source of the article as distributed in the arXiv e-print for this exact version is bundled unchanged as \texttt{sources/201759/source0.tex}.
\begin{definition}\label{def:admissible-bregman}
Let $\E$ be reflexive. A function $f:\E\to\R$ is called an \emph{admissible Bregman generator} if
\begin{enumerate}
 \item $f$ is Legendre and strongly coercive;
 \item $f$ is bounded and uniformly Fr\'echet differentiable on bounded subsets of $\E$;
 \item $f$ is uniformly convex on bounded subsets of $\E$;
 \item for every $x\in\E$ and $r>0$, the right Bregman sublevel set
 \[
  \{y\in\E:D_f(x,y)\leq r\},
 \]
 is bounded.
\end{enumerate}
\end{definition}

 \begin{equation}\label{eq:bregman-consistency}
  D_f(x_n,y_n)\to0
  \quad\Longleftrightarrow\quad
  \|x_n-y_n\|\to0.
 \end{equation}

\begin{equation}\label{eq:V-identity}
 V_f(x,x^*)=D_f\bigl(x,\nabla f^*(x^*)\bigr),
\end{equation}

\begin{equation}\label{eq:bregman-error}
 D_f(p,T_kx)\leq D_f(p,x)+\varepsilon_k,
 \qquad x\in\E,\ p\in F,\ k\geq1.
\end{equation}

\begin{theorem}\label{thm:bregman}
Let $\E$ be a reflexive real Banach space and let $f$ be an admissible Bregman generator. Let $\{T_k\}_{k\geq1}$ be mappings from $\E$ into itself, and let $\Gamma$ be a family of closed mappings such that
\[
 \bigcap_{k=1}^{\infty}\Fix(T_k)=F(\Gamma)\neq\varnothing.
\]
Assume that $F(\Gamma)$ is closed and convex, that $\{T_k\}$ satisfies the NST-condition with $\Gamma$, and that \eqref{eq:bregman-error} holds with $F=F(\Gamma)$. Let $0\leq\alpha_k\leq\overline\alpha<1$ for every $k$, and let $\{\beta_k\}\subset\R$ be bounded. Given $v_0\in\E$, set $v_1=v_0$ and $\C_1=\E$. Define
\begin{equation}\label{eq:bregman-algorithm}
\begin{aligned}
 w_k&=v_k+\beta_k(v_k-v_{k-1}),\\
 y_k&=\nabla f^*\!\left(\alpha_k\nabla f(w_k)
 +(1-\alpha_k)\nabla f(T_kw_k)\right),\\
 \C_{k+1}
 &=\bigl\{v\in\C_k:
 D_f(v,y_k)\leq D_f(v,w_k)+(1-\alpha_k)\varepsilon_k\bigr\},\\
 v_{k+1}&=P_{\C_{k+1}}^f v_0.
\end{aligned}
\end{equation}
Then every $\C_k$ is nonempty, closed, and convex, the iteration is well defined, and
\[
 v_k\longrightarrow P_{F(\Gamma)}^f v_0
\]
strongly.
\end{theorem}

\begin{proof}
Set $F:=F(\Gamma)$ and $\delta_k=(1-\alpha_k)\varepsilon_k$.

\smallskip
\noindent\emph{Step 1: The Sets $\C_k$ Are Closed and Convex and Contain $F$.}
The assertion holds at $k=1$ because $\C_1=\E$ is nonempty, closed, and convex, and $F\subset\C_1$. Suppose inductively that $\C_k$ is nonempty, closed, and convex and that $F\subset\C_k$. For $v\in\C_k$, expansion of the two Bregman distances gives
\begin{align*}
 D_f(v,y_k)-D_f(v,w_k)
 &=\langle v,\nabla f(w_k)-\nabla f(y_k)\rangle\\
 &\quad+f(w_k)-f(y_k)
 +\langle y_k,\nabla f(y_k)\rangle
 -\langle w_k,\nabla f(w_k)\rangle.
\end{align*}
Hence, setting
\[
 d_k=\delta_k+f(y_k)-f(w_k)
 +\langle w_k,\nabla f(w_k)\rangle
 -\langle y_k,\nabla f(y_k)\rangle,
\]
we obtain
\begin{equation}\label{eq:bregman-halfspace}
 \C_{k+1}
 =\{v\in\C_k:\langle v,\nabla f(w_k)-\nabla f(y_k)\rangle\leq d_k\}.
\end{equation}
Thus $\C_{k+1}$ is closed and convex.

Since $\nabla f^*=(\nabla f)^{-1}$, the definition of $y_k$ gives
\[
 \nabla f(y_k)
 =\alpha_k\nabla f(w_k)+(1-\alpha_k)\nabla f(T_kw_k).
\]
Let $p\in F$. Using \eqref{eq:V-identity}, convexity of $V_f(p,\cdot)$, and \eqref{eq:bregman-error},
\begin{align*}
 D_f(p,y_k)
 &\leq \alpha_kD_f(p,w_k)+(1-\alpha_k)D_f(p,T_kw_k)\\
 &\leq D_f(p,w_k)+\delta_k.
\end{align*}
Therefore $p\in\C_{k+1}$. Induction shows that $\C_k$ is nonempty, closed, and convex for every $k\geq1$.

\smallskip
\noindent\emph{Step 2: Boundedness and the Nested-Set Estimate.}
Fix $p\in F$. The Bregman projection inequality gives
\begin{equation}\label{eq:bregman-bounded}
 D_f(p,v_k)+D_f(v_k,v_0)\leq D_f(p,v_0).
\end{equation}
The map $z\mapsto D_f(z,v_0)=f(z)-\langle z,\nabla f(v_0)\rangle+f^*(\nabla f(v_0))$ is coercive. Hence the boundedness of $\{D_f(v_k,v_0)\}$ in \eqref{eq:bregman-bounded} implies that $\{v_k\}$ is bounded. Since $v_{k+1}\in\C_{k+1}\subset\C_k$,
\[
 D_f(v_{k+1},v_k)+D_f(v_k,v_0)\leq D_f(v_{k+1},v_0).
\]
Thus $\{D_f(v_k,v_0)\}$ is nondecreasing and bounded. If $m>n$, then $v_m\in\C_n$, and
\begin{equation}\label{eq:bregman-nested}
 D_f(v_m,v_n)
 \leq D_f(v_m,v_0)-D_f(v_n,v_0).
\end{equation}

\smallskip
\noindent\emph{Step 3: Convergence and Residual Decay.}
Let $a_k=D_f(v_k,v_0)$. The preceding argument shows that $a_k$ converges. If $\{v_k\}$ were not Cauchy, there would exist $\rho>0$ and indices $m_j>n_j\to\infty$ such that $\|v_{m_j}-v_{n_j}\|\geq\rho$. By \eqref{eq:bregman-nested},
\[
 0\leq D_f(v_{m_j},v_{n_j})\leq a_{m_j}-a_{n_j}\longrightarrow0.
\]
Both subsequences are bounded, so \eqref{eq:bregman-consistency} gives $\|v_{m_j}-v_{n_j}\|\to0$, a contradiction. Thus $\{v_k\}$ is Cauchy. Let $v_k\to q$. For each fixed $n$, $v_k\in\C_k\subset\C_n$ for all $k\geq n$. Since $\C_n$ is closed, $q\in\C_n$, and hence $q\in\bigcap_n\C_n$. Boundedness of $\{\beta_k\}$ yields $w_k\to q$.

We next show that the auxiliary sequences are bounded. Since $\{w_k\}$ is bounded and both $f$ and $\nabla f$ are bounded on bounded sets, the identity
\[
 D_f(p,w_k)=f(p)-f(w_k)-\langle p-w_k,\nabla f(w_k)\rangle,
\]
shows that $\{D_f(p,w_k)\}$ is bounded. Equation \eqref{eq:bregman-error}, together with boundedness of the convergent sequence $\{\varepsilon_k\}$, then shows that $\{D_f(p,T_kw_k)\}$ is bounded. The right-sublevel assumption in Definition~\ref{def:admissible-bregman} implies that $\{T_kw_k\}$ is bounded. Consequently, $\{\nabla f(w_k)\}$ and $\{\nabla f(T_kw_k)\}$ are bounded in $\E^*$. Their convex combinations are bounded, and boundedness of $\nabla f^*$ on bounded subsets of $\E^*$ gives boundedness of $\{y_k\}$.

Since $v_{k+1}-w_k\to0$, \eqref{eq:bregman-consistency} gives $D_f(v_{k+1},w_k)\to0$. Since $v_{k+1}\in\C_{k+1}$, the inequality defining $\C_{k+1}$ in \eqref{eq:bregman-algorithm} gives
\[
 0\leq D_f(v_{k+1},y_k)
 \leq D_f(v_{k+1},w_k)+\delta_k\to0.
\]
Hence $\|v_{k+1}-y_k\|\to0$, and therefore $\|y_k-w_k\|\to0$. Uniform continuity of $\nabla f$ on bounded sets yields
\[
 \|\nabla f(y_k)-\nabla f(w_k)\|\to0.
\]
By the definition of $y_k$,
\[
 \nabla f(y_k)-\nabla f(w_k)
 =(1-\alpha_k)\bigl(\nabla f(T_kw_k)-\nabla f(w_k)\bigr).
\]
Since $1-\alpha_k\geq1-\overline\alpha>0$, it follows that
\[
 \|\nabla f(T_kw_k)-\nabla f(w_k)\|\to0.
\]
The two gradient sequences are bounded. Hence, by the uniform norm-to-norm continuity of $\nabla f^*$ on bounded subsets of $\E^*$,
\begin{align*}
 \|T_kw_k-w_k\|
 &=\bigl\|\nabla f^*(\nabla f(T_kw_k))
 -\nabla f^*(\nabla f(w_k))\bigr\|\longrightarrow0.
\end{align*}
Thus
\begin{equation}\label{eq:bregman-residual}
 \|T_kw_k-w_k\|\to0.
\end{equation}

\smallskip
\noindent\emph{Step 4: Membership in $F$.}
By the NST-condition and \eqref{eq:bregman-residual},
\[
 \|Tw_k-w_k\|\to0,
 \qquad T\in\Gamma.
\]
Since $w_k\to q$ and $\|Tw_k-w_k\|\to0$, one has $Tw_k\to q$. Closedness of the graph of $T$ therefore gives $Tq=q$. Thus $q\in F$.

\smallskip
\noindent\emph{Step 5: Identification of the Limit.}
From \eqref{eq:bregman-bounded},
\[
 D_f(v_k,v_0)\leq D_f(p,v_0),
 \qquad p\in F.
\]
Since $z\mapsto D_f(z,v_0)$ is continuous and $v_k\to q$,
\[
 D_f(q,v_0)=\lim_{k\to\infty}D_f(v_k,v_0)\leq D_f(p,v_0),\qquad p\in F.
\]
Thus $q$ is the unique minimizer of $z\mapsto D_f(z,v_0)$ over $F$, and hence $q=P_F^fv_0$.
\end{proof}

\end{document}
